\section{Boolean and connected candidates}
\label{sec:candidates}

Let \(A=(a_1,\ldots,a_r)\) be a tuple of distinct atoms, initially with
\(r\in\{2,3\}\).  Suppose its join \(j(A)=\bigvee_i a_i\) exists.  The
subset-join map is
\begin{equation}
 \beta_A:2^{[r]}\longrightarrow[\bot,j(A)],
 \qquad S\longmapsto\bigvee_{i\in S}a_i,
 \label{eq:beta}
\end{equation}
with empty join \(\bot\).

\begin{definition}[Normalized Boolean-join weight]
\begin{equation}
 \KBool_r(A)
 :=\one_{\{\beta_A\text{ is a pointed order isomorphism}\}}
   (-1)^r\mu_X(\bot,j(A)).
 \label{eq:boolean-weight}
\end{equation}
If the join is undefined, the value is zero.
\end{definition}

On a Boolean lattice \(B_r\),
\(
 \mu(\hat0,\hat1)=(-1)^r
\), so every accepted tuple has \(\KBool_r=1\).

\begin{proposition}[Naturality and compatible cutoffs]
\label{prop:boolean-natural}
The weight \eqref{eq:boolean-weight} is invariant under admissible pointed
isomorphisms.  Its value on an old tuple is unchanged under a compatible
active-cutoff enlargement evaluated in the same ambient source.
\end{proposition}
\begin{proof}
An isomorphism carries atoms to atoms, preserves every subset join, and
restricts to an isomorphism of the two pointed join intervals.  It therefore
preserves the Boolean certificate and the top M{\"o}bius value.  An active
cutoff changes only which already compiled atoms enter a finite sum; it does
not replace the ambient interval of an old tuple.
\end{proof}

\begin{proposition}[Conditional uniqueness and scheme freedom]
\label{prop:conditional-unique}
Equation~\eqref{eq:boolean-weight} is the unique \(\{0,1\}\)-valued tuple
weight supported exactly on Boolean tuple-join intervals and normalized to
one there.  Isomorphism naturality alone does not select it.
\end{proposition}
\begin{proof}
The support and normalization axioms fix both possible values.  The M{\"o}bius
factor realizes the normalization because the Boolean top value is
\((-1)^r\).  Conversely, let \(\mathfrak I_r\) be the set of pointed
isomorphism classes of tuple-join intervals.  Every function
\(f:\mathfrak I_r\to\C\) yields the natural weight
\(
 f([\bot,j(A)]_{\rm pt})
\).  Global filtered-source data permit still more choices.
\end{proof}

The proposition separates a useful exact formula from an overclaim.
“Natural” does not mean “canonical”; Boolean support, range, and unit
normalization are scheme axioms.

\subsection{The connected alternative}

Fix a moment \(m(A_S)\) for each nonempty subset.  Its partition-lattice
connected transform is
\begin{equation}
 \kappa_r(A)=
 \sum_{\pi\in\Pi_r}(-1)^{|\pi|-1}(|\pi|-1)!
 \prod_{B\in\pi}m(A_B).
 \label{eq:cumulant}
\end{equation}
The coefficients in \eqref{eq:cumulant} are classical
\citep{Speed1983}; the source does not thereby choose \(m\).

\begin{lemma}[Factorization cancellation]
\label{lem:factorization-cancellation}
If
\(
 m(A_S)=\prod_{i\in S}u_i
\)
for every nonempty \(S\), then \(\kappa_r(A)=0\) for every \(r\ge2\).
\end{lemma}
\begin{proof}
Every partition product in \eqref{eq:cumulant} equals
\(\prod_i u_i\).  The remaining coefficient is
\[
 c_r=\sum_{\pi\in\Pi_r}(-1)^{|\pi|-1}(|\pi|-1)!.
\]
Its exponential generating function is
\(
 \sum_{r\ge1}c_r z^r/r!=\log(e^z)=z
\), so \(c_r=0\) for \(r\ge2\).
\end{proof}

For the Boolean moment on a free commutative source, every nonempty subset
moment is one.  Thus
\begin{equation}
 \kappa_2=1-1=0,
 \qquad
 \kappa_3=1-1-1-1+2=0.
 \label{eq:low-cumulants}
\end{equation}
The connected transform detects failure of multiplicative factorization;
it removes the desired free baseline signal.

\subsection{The frozen five-predicate refinement}

The exact suite strengthened the Boolean support before execution.  For a
pair or triple \(A\), define \(\chi_r(A)=1\) precisely when all five
conditions hold:
\begin{enumerate}
\item the join \(j(A)\) is unique;
\item \(\beta_A\) identifies its pointed interval with \(B_r\);
\item \(\mu(\bot,j(A))=(-1)^r\);
\item the roof is multiplicative,
  \(w(j(A))=\prod_{a\in A}w(a)\); and
\item iterated binary joins own \(j(A)\) associatively.
\end{enumerate}
All five predicates are source-resident and natural under the admitted
decorated isomorphisms.  The 31 nonempty predicate masks were enumerated as
a robustness audit; none was selected after observing controls.

The rank-three connected contraction used by that suite is
\begin{equation}
 \Theta_3
 =2\sum_{a<b<c}\chi_3(a,b,c)
 \frac{G_{ab}G_{bc}G_{ca}}{w(a)w(b)w(c)}.
 \label{eq:theta3}
\end{equation}
It is important not to conflate \(\Theta_3\) with the partition-lattice
cumulant \eqref{eq:cumulant}.  It is a separately filtered triangle
contraction, not the full \(\operatorname{Tr}\cB_s^6\) and not a chiral
determinant coefficient.  It can be nonzero on a factorizing baseline.

The final exact audit finds \(\chi_3\) and \(\Theta_3\) nonzero at all three
integer cutoffs and exactly zero on the four finite non-UFD fixtures.  This
is a genuine finite-fixture GO.  It is not arithmetic selectivity: every
pair, triple, coefficient, and predicate mask is reproduced by the
transported free and polynomial-UFD clones.

\begin{figure}[t]
\centering
\begin{tikzpicture}[
  >=Latex,
  head/.style={draw=charcoal,thick,rounded corners=2pt,fill=softgray,
    text width=37mm,align=center,minimum height=10mm},
  cell/.style={draw,thick,rounded corners=2pt,text width=37mm,align=center,
    minimum height=13mm},
  note/.style={font=\small,align=center,text width=31mm},
  every node/.style={font=\small}
]
\node[head] (inthead) {integer divisibility};
\node[head,right=29mm of inthead] (clonehead) {formal UFD clone};

\node[cell,draw=warningamber,fill=warningamber!7,below=11mm of inthead] (ik)
  {Boolean tuple interval\\$\KBool_r=1$};
\node[cell,draw=warningamber,fill=warningamber!7,below=11mm of clonehead] (fk)
  {Boolean tuple interval\\$\KBool_r=1$};
\draw[<->,thick,warningamber] (ik) -- node[above,fill=white] {same} (fk);

\node[note,left=4mm of ik] (klab) {nonzero, but\\not selective};

\node[cell,draw=deepgreen,fill=deepgreen!7,below=12mm of ik] (ic)
  {factorizing moment\\$\kappa_r=0$ for $r\ge2$};
\node[cell,draw=deepgreen,fill=deepgreen!7,below=12mm of fk] (fc)
  {factorizing moment\\$\kappa_r=0$ for $r\ge2$};
\draw[<->,thick,deepgreen] (ic) -- node[above,fill=white] {same} (fc);

\node[note,left=4mm of ic] (clab) {selective zero, but\\zero on baseline};
\end{tikzpicture}

\caption{The two frozen candidates fail complementary gates.  The Boolean
weight is nonzero but survives on the formal UFD control.  The connected
cumulant of a factorizing moment vanishes on the control and on the integer
baseline.  Changing the moment adds scheme data but cannot break the
isomorphism proved next.}
\label{fig:candidate-dichotomy}
\end{figure}
